\item \textbf{Modelado del fenómeno de pegar y deslizar (stick-slip).}
Su pulgar rechina en un plato que acaba de lavar. Sus zapatos tenis rechinan en el piso del gimnasio. Las llantas de los autos rechinan con un arranque o frenado abrupto. Usted puede hacer cantar una copa al secar su dedo humedecido alrededor de su borde. Cuando el gis rechina en un pizarrón, usted puede ver que hace una hilera de rayas regularmente espaciadas. Como sugieren estos ejemplos, la vibración comúnmente resulta cuando la fricción actúa sobre un objeto elástico en movimiento. La oscilación no es un movimiento armónico simple, sino que se llama pegar y deslizar. Este problema modela un movimiento de pegar y deslizar.

Un bloque de masa $m$ se une a un soporte fijo mediante un resorte horizontal, con constante de fuerza $k$ y masa despreciable (figura P15.42). La ley de Hooke describe al resorte tanto en extensión como en compresión. El bloque descansa sobre una larga tabla horizontal, con la que tiene coeficiente de fricción estático $\mu_s$ y un coeficiente de fricción cinética $\mu_k$ menor. La tabla se mueve hacia la derecha con rapidez constante $v$. Suponga que el bloque pasa la mayor parte de su tiempo pegado a la tabla y en movimiento hacia la derecha, de modo que la rapidez $v$ es pequeña en comparación con la rapidez promedio que tiene el bloque mientras se desliza de regreso hacia la izquierda.
\begin{enumerate}
    \item Demuestre que la máxima extensión del resorte, desde su posición no estirada, es muy cercana a la que se conoce mediante $\mu_s mg/k$.
    \item Demuestre que el bloque oscila en torno a una posición de equilibrio en la que el resorte se estira en $\mu_k mg/k$.
    \item Grafique cualitativamente la posición del bloque con el tiempo.
    \item Demuestre que la amplitud del movimiento del bloque es
    $$ A = \frac{(\mu_s - \mu_k)mg}{k} $$
    \item Demuestre que el periodo del movimiento del bloque es
    $$ T = \frac{2(\mu_s - \mu_k)mg}{vk} + \pi \sqrt{\frac{m}{k}} $$
\end{enumerate}